La distribución de clientes en mora por grupo de edad muestra que la mayor concentración se encuentra en los rangos comprendidos entre
20 y 64 años
. En estos grupos, la cantidad de clientes con incumplimiento oscila aproximadamente entre
12 y 19 clientes
, lo que evidencia una distribución relativamente homogénea a lo largo de la mayor parte de la población analizada.
En contraste, los grupos de mayor edad presentan una menor cantidad de clientes en mora, siendo especialmente reducido el número de casos en los clientes de
65 a 69 años
. No obstante, esta diferencia puede estar influenciada por una menor cantidad de clientes pertenecientes a dichos rangos de edad, por lo que no constituye evidencia suficiente para concluir que la edad sea un factor determinante del riesgo de mora.
En conjunto, el análisis descriptivo no permite identificar un patrón suficientemente marcado que relacione el grupo de edad con el incumplimiento de los clientes. Por ello, se recomienda complementar este estudio mediante análisis estadísticos y modelos predictivos que incorporen variables adicionales, como el nivel de ingreso, la antigüedad laboral, el monto del préstamo y la relación cuota-ingreso, con el fin de identificar los factores que realmente inciden en el riesgo crediticio y fortalecer las estrategias de otorgamiento y seguimiento de préstamos.